\section{The Dirichlet half-plane and its quantitative algebra}
\label{sec:dirichlet}

The classical formulas in this section are recorded in
\cite[\S\S25.2,25.8]{DLMF}. Their proofs below use finite sums, explicit
majorants, and elementary exponential and logarithm constructions.
Complex parameters have certified primitive recursive names, together
with the indicated rational bounds. No search for a distance from a
boundary is implicit.

\subsection{A single tail estimate}

Put \(n^{-s}=\exp(-s\log n)\), using the real logarithm of a positive
integer, and write \(\sigma=\operatorname{Re}s\).
Fix \(D\geq1\) and suppose \(\sigma\geq1+1/D\).
Let \(Z_N(s)=\sum_{n=1}^N n^{-s}\).

\begin{lemma}[Dyadic tail]\label{lem:d-tail}
If \(N=2^{DL}\), \(L\geq0\), and \(M\geq N\), then
\[
 \sum_{n=N+1}^{M}|n^{-s}|\leq2D\,2^{-L}.                 \tag{D1}\label{eq:d-tail}
\]
These finite sums therefore define a primitive recursively approximable
complex number \(\zeta(s)\), with the same tail bound.
\end{lemma}
\begin{proof}
On the block \(2^a\leq n<2^{a+1}\), the sum is at most
\(2^a2^{-a(1+1/D)}=2^{-a/D}\).
Group \(D\) consecutive blocks, writing \(a=D(L+r)+b\), \(0\leq b<D\).
Their sum is at most \(D2^{-L-r}\). The finite geometric identity bounds
any finite collection of groups by \(2D2^{-L}\). Including extra terms
at the endpoints only increases this nonnegative majorant.

For precision \(2^{-k}\), take \(L=k+2D+4\), evaluate each of the
\(N=2^{DL}\) summands with error at most \(2^{-k-3}/N\), and round the
sum to a dyadic number with error at most \(2^{-k-2}\).
Since \(2D\leq2^{2D}\), the truncation error is at most \(2^{-k-4}\).
All lengths and arithmetic accuracies are given in advance by primitive
recursive functions. The finite tail inequality supplies the common
modulus for these approximations.
\end{proof}

Finite triangle inequalities and termwise conjugation yield
\[
 |\zeta(s)|\leq\zeta(\sigma)\leq1+2D,\qquad
 \zeta(\overline s)=\overline{\zeta(s)}.                 \tag{D2}\label{eq:d-bound}
\]
Conversely, if \(\sigma\leq1\), each block
\(2^{a-1}<n\leq2^a\) contributes at least \(1/2\) to the sum of absolute
values. Thus absolute convergence holds exactly when \(\sigma>1\).
This statement concerns absolute convergence, and makes no inference
from its failure to failure of ordinary convergence.

\subsection{Derivatives and a local analytic expansion}

\begin{lemma}[Logarithmic weights]\label{lem:d-derivative-tail}
For \(j\geq0\), at the cutoff \(N=2^{2DL}\),
\[
 \sum_{n>N}(\log n)^j n^{-\sigma}
       \leq4D\,j!(2D)^j2^{-L}.                         \tag{D3}\label{eq:d-derivative-tail}
\]
Moreover
\(\zeta^{(j)}(s)=\sum_{n\geq1}(-\log n)^j n^{-s}\).
\end{lemma}
\begin{proof}
The positive exponential series gives \(y^j\leq j!e^y\).
With \(y=(\log n)/(2D)\),
\[
 (\log n)^j n^{-\sigma}
       \leq j!(2D)^j n^{-1-1/(2D)}.
\]
Apply Lemma~\ref{lem:d-tail} with \(2D\) in place of \(D\).

Here is also a quantitative justification of differentiability.
The finite exponential remainder obeys
\[
 \left|e^z-\sum_{j=0}^r z^j/j!\right|
       \leq |z|^{r+1}e^{|z|}/(r+1)! .
\]
For \(|h|\leq1/(8D)\), put \(z=-h\log n\) and use
\((\log n)^{r+1}\leq(r+1)!(4D)^{r+1}n^{1/(4D)}\).
The remaining nonnegative series is bounded by
\(\sum n^{-1-1/(4D)}\leq1+8D\). Consequently, with
\(c_j(s)=\sum_{n\geq1}(-\log n)^j n^{-s}/j!\),
\[
 \left|\zeta(s+h)-\sum_{j=0}^{r}c_j(s)h^j\right|
       \leq(1+8D)(4D|h|)^{r+1}.                       \tag{D4}\label{eq:d-taylor}
\]
Every passage from finite sums is controlled by \eqref{eq:d-derivative-tail}.
The right side is at most \((1+8D)2^{-r-1}\) on this disk.
Thus this is a local power series with an explicit radius and remainder,
not an appeal to an abstract uniform-limit theorem.
Finite Taylor coefficient identities, or the same estimate for the
differentiated series, establish the asserted derivatives.
\end{proof}

\subsection{Finite convolution and arithmetic coefficients}

Suppose two series have absolute masses at most \(A_0,B_0\) and absolute
tail bounds \(A(M),B(M)\). For their Dirichlet convolution
\(c(n)=\sum_{d\mid n}a(d)b(n/d)\), the triangular sum through \(M^2\)
differs from the product of the full sums by at most
\[
                         A(M)B_0+A_0B(M).             \tag{D5}\label{eq:d-convolution}
\]
Indeed, expand finite rectangular sums. If a pair \((u,v)\) is omitted
by \(uv\leq M^2\), then \(u>M\) or \(v>M\). Sum these two nonnegative
majorants, allowing double counting, and enlarge the outer rectangle
using the given tails. This proves the rearrangement and its modulus.
For logarithmic or polynomial weights, use
\eqref{eq:d-derivative-tail} or the corresponding shifted half-plane.

All coefficient functions below are primitive recursive. Primality and
divisibility are decided by trial division bounded by the input.
Factorizations are finite lists of primes and exponents, obtained by
bounded division. The identities below are finite product expansions
and bounded divisor sums; their inductions can be expressed with
quantifier-free invariants.

Define \(\mu(1)=1\), let \(\mu(n)=0\) if a prime square divides \(n\),
and otherwise put \(\mu(n)=(-1)^r\), where \(r\) is its number of prime
divisors. For \(n>1\), with distinct prime divisors \(p_1,\ldots,p_r\),
\[
 \sum_{d\mid n}\mu(d)
   =\sum_{S\subseteq\{1,\ldots,r\}}(-1)^{|S|}
   =(1-1)^r=0.                                       \tag{D6}\label{eq:d-mu}
\]
For \(n=1\), the sum is \(1\). Since \(|\mu(n)|\leq1\),
\(M(s)=\sum_{n\geq1}\mu(n)n^{-s}\) has tail \eqref{eq:d-tail}.
Finite convolution now gives
\[
                         \zeta(s)M(s)=1.              \tag{D7}\label{eq:d-reciprocal}
\]
As \(|M(s)|\leq1+2D\), we get actual apartness:
\[
                \boxed{|\zeta(s)|\geq(1+2D)^{-1}.}    \tag{D8}\label{eq:d-nonzero}
\]
This supplies the data needed to compute the reciprocal.

\subsection{Euler product with a prime-cutoff modulus}

For \(y\geq2\), let
\(P_y(s)=\prod_{p\leq y,\ p\ {\rm prime}}(1-p^{-s})^{-1}\).
Each denominator satisfies \(|1-p^{-s}|\geq1/2\).

\begin{theorem}[Euler product]\label{thm:d-euler-product}
If \(y\geq2^{DL}\), then
\[
                         |\zeta(s)-P_y(s)|\leq2D2^{-L}.
                                                               \tag{D9}\label{eq:d-product-tail}
\]
\end{theorem}
\begin{proof}
List the primes at most \(y\) as \(p_1,\ldots,p_m\). For a finite \(K\),
\[
 \prod_{i=1}^m\sum_{e=0}^Kp_i^{-es}
 =\sum_{0\leq e_1,\ldots,e_m\leq K}
                    (p_1^{e_1}\cdots p_m^{e_m})^{-s}.
\]
Unique factorization says that no integer on the right occurs twice.
Each geometric factor differs from its full value by at most \(2^{-K}\);
all factors have absolute value at most \(2\).
The finite telescoping product identity bounds the error by
\(m2^{m-1}2^{-K}\).

Fix a further finite cutoff \(T\). Once \(K\geq T\), all integers at most
\(T\) whose prime factors are at most \(y\) occur in the displayed
product. Remaining terms have absolute sum at most the tail of
\(\sum n^{-\sigma}\) beyond \(T\).
These two explicit bounds prove
\[
 P_y(s)=\sum_{\substack{n\geq1\\p\mid n\Rightarrow p\leq y}}n^{-s}.
\]
In particular \(|P_y(s)|\leq\zeta(\sigma)\).
Every \(n\leq y\) has only prime factors at most \(y\); hence the
difference from \(\zeta\) is a subseries indexed by \(n>y\), and
\eqref{eq:d-tail} proves the result.
\end{proof}

The reciprocal product has the finite expansion
\[
 \prod_{p\leq y}(1-p^{-s})
       =\sum_{d\mid\prod_{p\leq y}p}\mu(d)d^{-s}.
\]
Its difference from \(M(s)\) is supported on squarefree integers having
a prime factor greater than \(y\); it has the same tail bound.

\subsection{Euler's logarithm and the \(3\)--\(4\)--\(1\) inequality}

Let \(b(p^a)=1/a\) for prime powers, \(a\geq1\), and \(b(n)=0\) otherwise.
Prime-power recognition is a bounded calculation. The absolutely
convergent series
\[
 \mathcal L(s)=\sum_{n\geq2}b(n)n^{-s}
             =\sum_p\sum_{a\geq1}p^{-as}/a             \tag{D10}\label{eq:d-log-series}
\]
has tail \eqref{eq:d-tail}, and
\[
                         e^{\mathcal L(s)}=\zeta(s).  \tag{D11}\label{eq:d-log-zeta}
\]
This is the logarithm tending to \(0\) along the positive real axis.
It is not asserted to be the principal complex logarithm everywhere.

For completeness, put \(\ell(z)=\sum_{a\geq1}z^a/a\), \(|z|\leq1/2\).
Its tail after \(K\) is at most \(2|z|^{K+1}\).
The finite-coefficient power-series calculation for
\(F(z)=\exp(\ell(z))\) gives \((1-z)F'(z)=F(z)\), \(F(0)=1\).
Writing \(F(z)=\sum c_mz^m\), coefficient comparison gives
\((m+1)c_{m+1}=(m+1)c_m\), so \(c_m=1\) by induction.
The geometric and exponential tails justify evaluation and prove
\(e^{\ell(z)}=(1-z)^{-1}\).
Apply this to each \(z=p^{-s}\) in a finite prime product.
Omitted prime powers have indices greater than the prime cutoff.
Moreover, on \(|u|,|v|\leq2D\),
\[
                    |e^u-e^v|\leq3^{2D}|u-v|.
\]
This follows from the exponential series along the line segment, using
\(e\leq3\). Combine this modulus with \eqref{eq:d-product-tail} to
obtain \eqref{eq:d-log-zeta}.

For real \(\sigma>1,t\), put \(\theta=at\log p\). Then
\[
\begin{split}
3\mathcal L(\sigma)+4\operatorname{Re}\mathcal L(\sigma+it)
             +\operatorname{Re}\mathcal L(\sigma+2it)
 &=\sum_{p,a}\frac{p^{-a\sigma}}a
                  (3+4\cos\theta+\cos2\theta)\\
 &=\sum_{p,a}\frac{2p^{-a\sigma}}a(1+\cos\theta)^2\geq0 .
\end{split}
\]
Finite partial sums are nonnegative, and the absolute tail preserves the
inequality. Exponentiation therefore gives
\[
 \boxed{1\leq\zeta(\sigma)^3|\zeta(\sigma+it)|^4
                                  |\zeta(\sigma+2it)|.}
                                                               \tag{D12}\label{eq:d-341}
\]
This is the inequality used in the separate boundary-line argument.

\subsection{The von Mangoldt series}

Set \(\Lambda(p^a)=\log p\) for prime powers, and \(0\) otherwise.
Then \(0\leq\Lambda(n)\leq\log n\).
The factorization of \(n\) gives the finite identity
\[
                    \sum_{d\mid n}\Lambda(d)=\log n,  \tag{D13}\label{eq:d-mangoldt-conv}
\]
because a prime of exponent \(e\) contributes \(e\log p\).
Convolution and the derivative tail give
\(\zeta(s)\sum\Lambda(n)n^{-s}=-\zeta'(s)\).
Using \eqref{eq:d-nonzero},
\[
 -\frac{\zeta'(s)}{\zeta(s)}
   =\sum_{n\geq1}\frac{\Lambda(n)}{n^s}
   =\sum_p\frac{\log p}{p^s-1}.                        \tag{D14}\label{eq:d-log-derivative}
\]
The last equality expands geometric series and is controlled by
\(\sum(\log n)n^{-\sigma}\).
Also \(\mathcal L'(s)=-\sum\Lambda(n)n^{-s}\), since
\(b(p^a)\log(p^a)=\log p\).

\subsection{Divisor functions, totients, squarefree integers, and Liouville's function}

With a supplied rational gap from each boundary, one has
\[
\begin{array}{c|c|c}
\text{coefficient}&\text{Dirichlet series}&\text{domain}\\ \hline
\tau(n)=\sum_{d\mid n}1
 &\displaystyle\sum_{n\geq1}\tau(n)n^{-s}=\zeta(s)^2&\sigma>1\\[3pt]
\sigma_r(n)=\sum_{d\mid n}d^r,\ r\in\mathbb N
 &\displaystyle\sum_{n\geq1}\sigma_r(n)n^{-s}=\zeta(s)\zeta(s-r)&\sigma>r+1\\[3pt]
\varphi(n)=\#\{1\leq a\leq n:(a,n)=1\}
 &\displaystyle\sum_{n\geq1}\varphi(n)n^{-s}=\zeta(s-1)/\zeta(s)&\sigma>2\\[3pt]
\mu(n)^2
 &\displaystyle\sum_{n\geq1}\mu(n)^2n^{-s}=\zeta(s)/\zeta(2s)&\sigma>1\\[3pt]
\lambda(n)=(-1)^{\sum_p e_p(n)}
 &\displaystyle\sum_{n\geq1}\lambda(n)n^{-s}=\zeta(2s)/\zeta(s)&\sigma>1 .
\end{array}                                                   \tag{D15}\label{eq:d-arithmetic-series}
\]
For \(\tau,\sigma_r\), their defining finite divisor sums are the
required convolutions. For the totient, use the M\"obius identity for
the gcd and exchange finite sums:
\[
 \varphi(n)=\sum_{a=1}^n\sum_{d\mid(a,n)}\mu(d)
           =\sum_{d\mid n}\mu(d)n/d.
\]
There are exactly \(n/d\) multiples of \(d\) among \(1,\ldots,n\).
Also \(\varphi(n)\leq n\), supplying the stated convergence domain.
For the squarefree coefficient,
\[
                         \mu(n)^2=\sum_{d^2\mid n}\mu(d).
\]
Indeed the right side factors into \(1-1\) for each prime whose exponent
in \(n\) is at least \(2\), and equals \(1\) if there is no such prime.
Convolution with the coefficient supported on squares gives
\(\sum\mu(n)^2n^{-s}=\zeta(s)M(2s)\).
Finally,
\[
 \sum_{d\mid n}\lambda(d)
     =\prod_{p\mid n}\sum_{j=0}^{e_p(n)}(-1)^j
     =
     \begin{cases}1,&n\text{ is a square},\\0,&\text{otherwise}.\end{cases}
\]
Thus \(\zeta(s)\sum\lambda(n)n^{-s}=\zeta(2s)\).
All infinite operations use \eqref{eq:d-convolution}; every division
uses the quantitative apartness \eqref{eq:d-nonzero}.

\subsection{Odd terms and the alternating series}

Deleting multiples of \(2\) in finite sums gives
\[
 \sum_{j\geq0}(2j+1)^{-s}=(1-2^{-s})\zeta(s),
              \qquad \sigma>1.                       \tag{D16}\label{eq:d-odd}
\]
Here \(1-2^{-s}\) is apart from zero.
For the alternating series, assume instead
\(\sigma\geq1/D,\ |s|\leq B,\ D,B\geq1\), and put
\(\eta_N(s)=\sum_{n=1}^N(-1)^{n-1}n^{-s}\).
Finite summation by parts, with partial sums of the alternating signs
bounded by \(1\), gives
\[
 \left|\sum_{n=N+1}^M(-1)^{n-1}n^{-s}\right|
 \leq M^{-\sigma}+|s|\int_{N+1}^{M}x^{-\sigma-1}\,dx
 \leq(1+BD)N^{-\sigma}.                               \tag{D17}\label{eq:d-eta-tail}
\]
The bound on consecutive differences used here is
\[
 |n^{-s}-(n+1)^{-s}|
     \leq |s|\int_n^{n+1}x^{-\sigma-1}\,dx.
\]
This is the finite-interval fundamental theorem for the explicitly
differentiated \(x^{-s}\), equivalently its rational Riemann-sum
estimate. At \(N=2^{DL}\), the tail is at most \((1+BD)2^{-L}\).
It defines \(\eta(s)\) primitively recursively throughout \(\sigma>0\).

The exact identity
\(\eta_{2N}(s)=Z_{2N}(s)-2^{1-s}Z_N(s)\) gives
\[
                       \eta(s)=(1-2^{1-s})\zeta(s)    \tag{D18}\label{eq:d-eta-relation}
\]
first on \(\sigma>1\), and also for Euler--Maclaurin continuation on
\(\sigma>0\), \(s\) apart from \(1\). To prove the latter directly, use
that continuation's compatibility estimate
\[
 Z_N(s)=\zeta(s)+\frac{N^{1-s}}{1-s}+E_N(s),\qquad
 |E_N(s)|\leq
 \left(\frac12+\frac B{12}+\frac{B(B+1)}{12}\right)N^{-\sigma}.
                                                               \tag{D19}\label{eq:d-eta-em}
\]
The terms containing \(N^{1-s}\) cancel exactly. Both remaining errors
have the displayed primitive recursive decay.
No identity theorem is needed.

The quotient \(\eta(s)/(1-2^{1-s})\) requires an apartness witness for
its denominator. This factor also vanishes at
\(s=1+2\pi i k/\log2\), \(k\ne0\). At those points
\eqref{eq:d-eta-relation} says \(\eta(s)=0\); quotient notation alone
does not define \(\zeta(s)\). Euler--Maclaurin continuation does.

At \(s=1\), integrate the finite geometric identity:
\[
 \eta_N(1)=\int_0^1\frac{1-(-x)^N}{1+x}\,dx,\qquad
 |\eta_N(1)-\log2|\leq1/(N+1).
\]
Thus
\[
                              \eta(1)=\log2.          \tag{D20}\label{eq:d-eta-one}
\]
For a real \(x>0\), grouping consecutive pairs gives
\(\eta(x)\geq1-2^{-x}>0\).
More explicitly, if \(x\geq1/D\), Bernoulli's finite inequality
\((1-1/(2D))^D\geq1/2\) implies
\(2^{-1/D}\leq1-1/(2D)\), so \(\eta(x)\geq1/(2D)\).
For \(0<x<1\), the factor \(1-2^{1-x}\) is strictly negative.
With supplied rational gaps from both endpoints its inverse is
computable, and therefore
\[
                         \zeta(x)<0\quad(0<x<1).      \tag{D21}\label{eq:d-real-negative}
\]
In particular this real interval contains no zeta zero.

\subsection{Real monotonicity, convexity, and asymptotics}

For real \(x>1\), the derivative formulas give
\[
 (-1)^j\zeta^{(j)}(x)
       =\sum_{n\geq1}(\log n)^j n^{-x}>0\quad(j\geq1).
                                                               \tag{D22}\label{eq:d-complete-monotone}
\]
The \(n=2\) term gives the lower bound \((\log2)^j2^{-x}\).
Together with \(\zeta(x)-1>0\), this proves strict complete monotonicity
of \(\zeta-1\), and hence strict decrease and strict convexity of \(\zeta\).
Without derivatives, for \(h>0\),
\(\zeta(x)-\zeta(x+h)\geq2^{-x}(1-2^{-h})>0\).
Finite expansion of a square gives
\[
 \zeta(x)\zeta''(x)-\zeta'(x)^2
   =\frac12\sum_{m,n\geq1}m^{-x}n^{-x}(\log m-\log n)^2
   \geq2^{-x}(\log2)^2>0.                             \tag{D23}\label{eq:d-logconvex}
\]
The pairs \((1,2),(2,1)\) supply the lower bound; derivative tails
authorize the double sum. Division by \(\zeta(x)^2>0\) proves strict
logarithmic convexity.

Comparing the decreasing \(t^{-x}\) with its endpoint sums gives
\[
 \frac1{x-1}\leq\zeta(x)\leq1+\frac1{x-1},\qquad
 2^{-x}\leq\zeta(x)-1\leq2^{-x}+\frac{2^{1-x}}{x-1}.   \tag{D24}\label{eq:d-real-bounds}
\]
These comparisons are first made on finite intervals; their omitted
integral is \(T^{1-x}/(x-1)\), with explicit dyadic decay.
In particular \(1\leq\delta\zeta(1+\delta)\leq1+\delta\), \(\delta>0\).
For \(x\geq2\) one also has the fully discrete bound
\[
                2^{-x}\leq\zeta(x)-1\leq3\cdot2^{-x}. \tag{D25}\label{eq:d-three}
\]
Indeed \(\sum_{n\geq2}n^{-2}\leq1/4+
\sum_{n\geq3}1/(n(n-1))=3/4\); multiply by \(2^{2-x}\).
Hence \(\zeta(x)\to1\) with an explicit exponential modulus.
The \(n\geq3\) terms, bounded by their first term plus
\(\int_3^\infty t^{-x}dt\), give
\[
 0\leq2^x(\zeta(x)-1)-1
       \leq(2/3)^x\left(1+\frac3{x-1}\right).           \tag{D26}\label{eq:d-leading-asymptotic}
\]
This proves effectively that \(2^x(\zeta(x)-1)\to1\).

\subsection{Sums of zeta values and a generating function}

\begin{theorem}\label{thm:d-zeta-sums}
With primitive recursive moduli,
\[
\begin{aligned}
 \sum_{k=2}^{\infty}(\zeta(k)-1)&=1,&
 \sum_{j=1}^{\infty}(\zeta(2j)-1)&=\tfrac34,\\
 \sum_{j=1}^{\infty}(\zeta(2j+1)-1)&=\tfrac14,&
 \sum_{k=2}^{\infty}(-1)^k(\zeta(k)-1)&=\tfrac12.
\end{aligned}                                                   \tag{D27}\label{eq:d-zeta-sums}
\]
\end{theorem}
\begin{proof}
Finite geometric summation gives
\[
 \sum_{k=2}^K(\zeta(k)-1)
  =\sum_{n\geq2}\left(\frac1{n(n-1)}-\frac1{(n-1)n^K}\right).
\]
The first series telescopes to \(1\); the remainder is at most
\(\zeta(K)-1\leq3\cdot2^{-K}\).
The \(n\)-series has its own telescoping tail \(1/N\), so this is a
finite rearrangement with two supplied truncation parameters.

For the even exponents the geometric sum is
\(\sum_{n\geq2}1/(n^2-1)\), and
\[
 \frac1{n^2-1}=\frac12\left(\frac1{n-1}-\frac1{n+1}\right)
\]
telescopes to \(3/4\). The error after \(j=K\) is at most
\((3/4)4^{-K}\).
For odd exponents starting with \(3\), use
\[
 \frac1{n(n^2-1)}
   =\frac12\left(\frac1{n-1}+\frac1{n+1}-\frac2n\right)
\]
to obtain \(1/4\), with error at most \((1/4)4^{-K}\).
The alternating geometric sum is
\(\sum_{n\geq2}1/(n(n+1))=1/2\). Its error after \(k=K\) is at most
\(3\cdot2^{-K}\), by \eqref{eq:d-three} and a geometric tail.
\end{proof}

For \(|z|\leq\rho<2\), with rational \(\rho\) supplied,
\[
 \sum_{k=2}^{\infty}(\zeta(k)-1)z^k
       =\sum_{n=2}^{\infty}\frac{z^2}{n(n-z)}.          \tag{D28}\label{eq:d-generating}
\]
Finite geometric identities prove this with tails
\[
 \frac{3(\rho/2)^{K+1}}{1-\rho/2}\quad\hbox{and}\quad
 \frac{\rho^2}{(1-\rho/2)N}
\]
beyond \(k=K\) and \(n=N\), respectively.
The given positive rational margin \(2-\rho\) yields primitive
recursive decay, so no unlicensed infinite-sum interchange occurs.

\subsection{A finite arithmetic interpretation: coprime tuples}

For \(r\geq2,N\geq1\), let \(C_r(N)\) count tuples in
\(\{1,\ldots,N\}^r\) with gcd \(1\). M\"obius inversion of their gcd
gives the exact integer identity
\[
                  C_r(N)=\sum_{d=1}^{N}\mu(d)\lfloor N/d\rfloor^r.
                                                               \tag{D29}\label{eq:d-coprime-count}
\]
For \(0\leq v\leq u\), factoring \(u^r-v^r\) gives
\(u^r-v^r\leq r(u-v)u^{r-1}\).
Set \(u=1/d\), \(v=\lfloor N/d\rfloor/N\), and sum:
\[
 \left|\frac{C_r(N)}{N^r}-\sum_{d=1}^N\mu(d)d^{-r}\right|
       \leq\frac rN\sum_{d=1}^Nd^{-(r-1)}.
\]
Let \(\ell(N)\) be the binary length. Dyadic grouping gives
\(\sum_{d\leq N}1/d\leq1+\ell(N)\).
The omitted reciprocal series has tail at most
\(\sum_{d>N}d^{-2}\leq1/N\), by telescoping. Therefore
\[
 \boxed{\left|\frac{C_r(N)}{N^r}-\frac1{\zeta(r)}\right|
          \leq\frac{r(1+\ell(N))+1}{N}.}              \tag{D30}\label{eq:d-coprime-density}
\]
This is the precise finite meaning of the coprime-tuple density.
For instance, if \(b=\ell(r)\), then
\(N=2^{\,2(k+b+8)}\) makes the right side at most \(2^{-k}\);
the verification is an elementary exponential-versus-linear inequality.

\subsection{Arithmetic strength of the arguments}

All truncations are fixed combinations of addition, multiplication,
factorial, integer exponentiation, bounded factorization, finite sums,
and the elementary-function approximators. Series products reduce to
\eqref{eq:d-convolution}, differentiation to
\eqref{eq:d-derivative-tail}, and prime products to finite geometric
expansions. There is no unbounded search for a prime, a modulus, or a
nonzero value.

In the object theory the analytic notation abbreviates families of
finite rational error inequalities. The inductions concern finite
blocks, products, factorizations, and explicit recurrences, and can be
expressed as quantifier-free invariants on encoded data. The statements
are schematic in any fixed certified primitive recursive parameter name.
They do not postulate a universal primitive recursive evaluator for all
codes of primitive recursive programs.
